\documentclass[aspectratio=169]{beamer}

\usetheme{Madrid}
\usecolortheme{default}
\setbeamertemplate{navigation symbols}{}

\usepackage[T1]{fontenc}
\usepackage[utf8]{inputenc}
\usepackage{lmodern}
\usepackage{booktabs}
\usepackage{hyperref}

\title{LLM-Based Sign Language Translation}
\subtitle{From an Early RTMPoses Pipeline to a Better How2Sign Pose-to-English System}
\author{Your Name}
\institute{Your Department / College}
\date{\today}

\begin{document}

\begin{frame}
  \titlepage
\end{frame}

\begin{frame}{Introduction}
\begin{itemize}
  \item Sign language translation is an important task in accessibility and assistive AI.
  \item The aim is to convert sign language input into meaningful natural language text.
  \item In our project, we focused on training large language models for sign-to-English translation.
  \item Our work progressed in two stages: an early RTMPoses-based pipeline and a later improved How2Sign pipeline.
\end{itemize}
\end{frame}

\begin{frame}{Project Journey}
\begin{itemize}
  \item We did not reach the final pipeline in one step.
  \item Our first implementation used an older RTMPoses-based setup and produced very poor results.
  \item After studying the failure patterns, we redesigned the training setup.
  \item We then built a better decoder-only LLM pipeline on the How2Sign dataset.
\end{itemize}
\end{frame}

\begin{frame}{Problem Statement}
\begin{itemize}
  \item Given a sequence of pose features from sign language, generate the correct English sentence.
  \item The challenge is not only language generation, but also multimodal alignment.
  \item The model must connect body motion, hand movement, and temporal dynamics with text.
\end{itemize}
\end{frame}

\begin{frame}{Main Objective}
\begin{itemize}
  \item Explore whether large language models can be adapted for sign language translation.
  \item Build training pipelines that accept pose-based input and produce English output.
  \item Compare multiple LLM backbones.
  \item Study what failed in the first pipeline and what improved in the second one.
\end{itemize}
\end{frame}

\begin{frame}{Project Context}
\begin{itemize}
  \item Our work was inspired by the PoseStich-SLT repository.
  \item We used available support code and structure where useful.
  \item Our main work focused on the training pipelines, experiments, predictions, and analysis.
  \item The final better-performing setup used the How2Sign dataset.
\end{itemize}
\end{frame}

\begin{frame}{Two Phases of Our Work}
\begin{enumerate}
  \item First phase: old RTMPoses-based implementation
  \item Second phase: improved How2Sign pose-to-English LLM pipeline
\end{enumerate}
\vspace{0.5em}
\begin{itemize}
  \item The first phase helped us understand major failure modes.
  \item The second phase gave better results and a cleaner multimodal pipeline.
\end{itemize}
\end{frame}

\begin{frame}{Phase 1: Old RTMPoses-Based Setup}
\begin{itemize}
  \item The old implementation is stored in the folder \texttt{Sign-Language-Translation-old}.
  \item It used RTMPoses-based pose features stored as \texttt{.pkl} files.
  \item The pipeline was more generic and supported many model families.
  \item This was our first serious attempt, but the results were extremely poor.
\end{itemize}
\end{frame}

\begin{frame}{Old Pipeline Data Setup}
\begin{itemize}
  \item The old configs point to \texttt{rtmfeaturesisign} data paths.
  \item Pose files were loaded from a performance directory using \texttt{.pkl} inputs.
  \item The data loader used \texttt{266} input dimensions, corresponding to \texttt{133} joints with x and y coordinates.
  \item The sequence length was reduced using frame skipping with \texttt{step\_frames = 4}.
\end{itemize}
\end{frame}

\begin{frame}{Old Pipeline Input Representation}
\begin{itemize}
  \item Each frame in the old pipeline used \texttt{266}-dimensional RTMPoses-based features.
  \item The data loader handled variable pose lengths and padded them to a fixed limit.
  \item Spatial normalization was added to stabilize training.
  \item Gaussian noise augmentation was also supported.
\end{itemize}
\end{frame}

\begin{frame}{Old Pipeline Model Design}
\begin{itemize}
  \item The old code used a universal model factory.
  \item It supported both encoder-decoder models and decoder-only causal LMs.
  \item A temporal feature projector was added before the language model.
  \item LoRA-based fine-tuning was used for efficient adaptation.
\end{itemize}
\end{frame}

\begin{frame}{Old Pipeline Prompting Strategy}
\begin{itemize}
  \item The old model wrapper used prompt-style conditioning.
  \item One prompt format was: \texttt{User: Translate the following sign language kinematics into an English sentence.}
  \item The goal was to force the LLM into an instruction-following mode.
  \item In practice, this still did not produce grounded translations.
\end{itemize}
\end{frame}

\begin{frame}{Old Pipeline Training Setup}
\begin{itemize}
  \item Maximum frames were reduced to \texttt{75}.
  \item Batch size was typically \texttt{4}.
  \item Gradient accumulation was \texttt{4}.
  \item Learning rate for the LLM was \texttt{5e-5}.
  \item Projector learning rate was \texttt{1e-3}.
  \item Evaluation was run every \texttt{2} epochs.
  \item Beam search was turned off with \texttt{num\_beams = 1}.
\end{itemize}
\end{frame}

\begin{frame}{Old Models We Tried}
\begin{itemize}
  \item Qwen 1.5B
  \item Phi 1.5
  \item Llama 1B
  \item Gemma 2B
  \item Gemma 3 4B
  \item Gemma 7B
  \item Other configs were also prepared for larger models and seq2seq baselines
\end{itemize}
\end{frame}

\begin{frame}{Old Results Summary}
\begin{table}
\centering
\small
\begin{tabular}{lcc}
\toprule
Model & Logged Best BLEU & Best Epoch \\
\midrule
qwen\_1.5b & 0.001147 & 22 \\
phi\_1.5 & 0.001062 & 22 \\
llama\_1b & 0.000707 & 28 \\
gemma\_2b & 0.000403 & 4 \\
gemma\_3\_4b & 0.000336 & 28 \\
gemma\_7b & 0.000187 & 8 \\
\bottomrule
\end{tabular}
\end{table}
\end{frame}

\begin{frame}{What These Old Results Mean}
\begin{itemize}
  \item The BLEU values were extremely close to zero.
  \item This means the generated text had almost no meaningful overlap with the target translation.
  \item Even when the output looked fluent, it was usually unrelated to the reference.
  \item So the first pipeline clearly failed at semantic grounding.
\end{itemize}
\end{frame}

\begin{frame}{Qualitative Failure Pattern in the Old Pipeline}
\begin{itemize}
  \item Predictions often became long unrelated paragraphs.
  \item Some outputs looked like random textbook explanations or news-style stories.
  \item The model generated fluent language, but not faithful translation.
  \item This showed that the LLM prior was dominating the pose signal.
\end{itemize}
\end{frame}

\begin{frame}{Why the Old Pipeline Performed Poorly}
\begin{itemize}
  \item RTMPoses features alone were not enough to ground the translation well.
  \item The generic prompting strategy did not solve the multimodal alignment problem.
  \item The models learned to produce fluent generic text instead of correct target sentences.
  \item Better data alignment and a cleaner pose-to-text setup were needed.
\end{itemize}
\end{frame}

\begin{frame}{What We Learned from Phase 1}
\begin{itemize}
  \item Bigger models do not automatically solve sign translation.
  \item Prompting alone is not enough for strong multimodal grounding.
  \item Pose feature design and data compatibility matter a lot.
  \item We needed a better dataset setup and a more focused training pipeline.
\end{itemize}
\end{frame}

\begin{frame}{Phase 2: Improved How2Sign Pipeline}
\begin{itemize}
  \item After the first failure, we moved to a cleaner How2Sign-based setup.
  \item This became the main LLM training pipeline of our project.
  \item The final scripts are \texttt{train\_how2sign\_llm.py} and \texttt{train\_how2sign\_llm\_A100.py}.
  \item This pipeline gave clearly better results than the earlier RTMPoses setup.
\end{itemize}
\end{frame}

\begin{frame}{Dataset Used in the Improved Pipeline}
\begin{itemize}
  \item The improved setup uses the How2Sign dataset.
  \item How2Sign is a large-scale continuous American Sign Language dataset.
  \item It contains aligned sign sequences and English sentence annotations.
  \item In our pipeline, pose files were used as the model input and English text as the target output.
\end{itemize}
\end{frame}

\begin{frame}{Why We Shifted to How2Sign}
\begin{itemize}
  \item The earlier setup was not giving meaningful translations.
  \item We wanted a cleaner and better-aligned pose-to-English training setup.
  \item How2Sign provided a more direct continuous sign translation setting.
  \item This made it more suitable for our decoder-only LLM experiments.
\end{itemize}
\end{frame}

\begin{frame}{Improved Pipeline Input Representation}
\begin{itemize}
  \item In the improved setup, each frame is represented using \texttt{152} values.
  \item These come from \texttt{76} keypoints with x and y coordinates.
  \item The keypoints include left hand, right hand, body, and face information.
  \item This feature sequence becomes the pose input for the model.
\end{itemize}
\end{frame}

\begin{frame}{Improved Pose Preprocessing}
\begin{itemize}
  \item Pose files are read and converted into frame-wise vectors.
  \item Hand, body, and face features are combined into one sequence.
  \item Sequences are padded or truncated to a fixed maximum number of frames.
  \item Noise can be added during training for robustness.
\end{itemize}
\end{frame}

\begin{frame}{Improved LLM Input Format}
\begin{itemize}
  \item A special token \texttt{<pose>} marks the start of the pose segment.
  \item Placeholder slots are reserved for pose embeddings.
  \item A special token \texttt{<English>} marks the beginning of the text segment.
  \item The model is trained autoregressively to generate the English sentence.
\end{itemize}
\end{frame}

\begin{frame}{Improved Prompt Strategy}
\begin{itemize}
  \item The improved training script uses the prompt: \texttt{Translate ASL to English:}
  \item Prompt embeddings are inserted along with pose embeddings.
  \item This is cleaner and better suited to instruction-tuned LLMs.
  \item It fits more naturally into decoder-only generation.
\end{itemize}
\end{frame}

\begin{frame}{Improved Models We Trained}
\begin{itemize}
  \item Qwen2.5-1.5B
  \item Qwen2.5-7B-Instruct
  \item Meta-Llama-3.1-8B-Instruct
  \item gemma-3-12b-it
  \item Mistral-7B-Instruct-v0.3
  \item Phi-4 based run
\end{itemize}
\end{frame}

\begin{frame}{Core Idea of the Improved Pipeline}
\begin{itemize}
  \item Pose features cannot directly be understood by an LLM.
  \item So we added a learned feature projection module.
  \item The projection maps pose features into the hidden space of the selected language model.
  \item After this alignment, the model generates English conditioned on sign input.
\end{itemize}
\end{frame}

\begin{frame}{Feature Projection Module}
\begin{itemize}
  \item Spatial linear layer
  \item Temporal convolution layer
  \item GELU activation
  \item Final output projection to the hidden size of the language model
\end{itemize}
\vspace{0.5em}
This module is the bridge between pose representation and language generation.
\end{frame}

\begin{frame}{Improved Training Method}
\begin{itemize}
  \item Fine-tuning was done using LoRA.
  \item The model was trained in bfloat16 precision.
  \item Flash Attention 2 was enabled in the main script.
  \item Gradient accumulation was used because of memory constraints.
\end{itemize}
\end{frame}

\begin{frame}{Important Hyperparameters in the Improved Setup}
\begin{table}
\centering
\small
\begin{tabular}{ll}
\toprule
Parameter & Value \\
\midrule
LLM learning rate & \texttt{5e-5} \\
Projection learning rate & \texttt{1e-3} \\
Maximum frames & \texttt{600} \\
Maximum decoder length & \texttt{128} \\
Batch size & \texttt{1} \\
Gradient accumulation & \texttt{16} \\
Epochs & \texttt{100} \\
LoRA rank & \texttt{16} \\
LoRA alpha & \texttt{32} \\
\bottomrule
\end{tabular}
\end{table}
\end{frame}

\begin{frame}{Outputs Saved in the Improved Pipeline}
\begin{itemize}
  \item Model checkpoints
  \item Validation predictions
  \item Test predictions
  \item Metrics saved in Excel format
\end{itemize}
\vspace{0.5em}
This allowed direct comparison across different LLM backbones and epochs.
\end{frame}

\begin{frame}{Evaluation Metrics}
\begin{itemize}
  \item BLEU-1
  \item BLEU-4
  \item SacreBLEU
  \item ROUGE-L
  \item Validation loss
  \item Test loss
\end{itemize}
\end{frame}

\begin{frame}{Best Validation Results in the Improved Setup}
\begin{table}
\centering
\small
\begin{tabular}{lccc}
\toprule
Model & BLEU-4 & SacreBLEU & ROUGE-L \\
\midrule
Qwen2.5-1.5B & 1.68 & 1.78 & 11.95 \\
phi-4\_300 & 0.90 & 0.95 & 11.06 \\
Meta-Llama-3.1-8B-Instruct & 0.83 & 0.96 & 9.63 \\
Qwen2.5-7B-Instruct & 0.78 & 0.84 & 9.58 \\
Mistral-7B-Instruct-v0.3\_300 & 0.63 & 0.70 & 5.34 \\
gemma-3-12b-it & 0.17 & 0.21 & 6.94 \\
\bottomrule
\end{tabular}
\end{table}
\end{frame}

\begin{frame}{Comparison: Old vs Improved Pipeline}
\begin{itemize}
  \item Old RTMPoses pipeline best logged BLEU was around \texttt{0.0011}.
  \item Improved How2Sign pipeline reached a validation BLEU-4 of \texttt{1.68}.
  \item The improved setup was still challenging, but it was clearly much better.
  \item This shows that dataset setup and multimodal design made a major difference.
\end{itemize}
\end{frame}

\begin{frame}{Main Observation}
\begin{itemize}
  \item Qwen2.5-1.5B gave the best validation result among the improved runs.
  \item Larger models did not automatically outperform smaller ones.
  \item Better alignment mattered more than just more parameters.
  \item Sign language translation remains a difficult grounding problem.
\end{itemize}
\end{frame}

\begin{frame}{Qualitative Findings in the Improved Setup}
\begin{itemize}
  \item The models generated more sentence-like English than in the old setup.
  \item However, many outputs were still not fully faithful to the reference.
  \item Some predictions were fluent but off-topic.
  \item This means the system improved, but semantic grounding is still incomplete.
\end{itemize}
\end{frame}

\begin{frame}{Strengths of Our Final Work}
\begin{itemize}
  \item We explored the problem in multiple stages instead of relying on one experiment.
  \item We identified clear failure patterns in the old RTMPoses setup.
  \item We redesigned the pipeline and achieved better results on How2Sign.
  \item We generated predictions and metric summaries across multiple LLMs.
\end{itemize}
\end{frame}

\begin{frame}{Limitations}
\begin{itemize}
  \item Even the improved results are still modest.
  \item Fluent text does not always mean correct translation.
  \item Pose-to-language alignment remains hard.
  \item Better multimodal fusion and stronger supervision are still needed.
\end{itemize}
\end{frame}

\begin{frame}{Future Work}
\begin{itemize}
  \item Improve the pose encoder and projection module.
  \item Add stronger temporal modeling.
  \item Explore cross-attention based fusion methods.
  \item Use better data filtering and more detailed human evaluation.
  \item Continue testing more LLMs with a unified training setup.
\end{itemize}
\end{frame}

\begin{frame}{Future Work: Implementing SignBind-LLM}
\begin{itemize}
  \item Our next major direction is to implement ideas from the SignBind-LLM paper.
  \item SignBind-LLM is a multi-stage modality fusion framework for sign language translation.
  \item Instead of forcing one network to learn everything at once, it separates the problem into specialized parts.
  \item This is important because natural signing contains multiple information streams that are not equally easy to model.
\end{itemize}
\end{frame}

\begin{frame}{Why SignBind-LLM Matters for Our Project}
\begin{itemize}
  \item Our current pipeline mainly projects pose features into an LLM and directly generates English.
  \item This works to some extent, but our results show that semantic grounding is still weak.
  \item According to the SignBind-LLM paper, single-stream approaches often fail on high-speed fingerspelling and asynchronous non-manual facial cues.
  \item Those are exactly the kinds of signals that can carry important names, places, and technical terms.
\end{itemize}
\end{frame}

\begin{frame}{Core Idea of SignBind-LLM}
\begin{itemize}
  \item SignBind-LLM uses separate expert predictors for continuous signing, fingerspelling, and lipreading.
  \item Each expert first converts its own modality into a token sequence.
  \item These parallel token streams are then fused using a lightweight transformer.
  \item The fused representation is finally passed to a large language model for sentence generation.
\end{itemize}
\end{frame}

\begin{frame}{How We Plan to Use SignBind Ideas}
\begin{itemize}
  \item We want to move from a single-stream pose-to-LLM setup to a richer multi-stream architecture.
  \item One stream can focus on general continuous signing features.
  \item Another stream can focus on fingerspelling-like detailed hand motion.
  \item Another stream can capture lip and facial information as non-manual cues.
  \item After that, we can fuse these streams before final language generation.
\end{itemize}
\end{frame}

\begin{frame}{Expected Benefit of a SignBind-Style Pipeline}
\begin{itemize}
  \item Better handling of temporal misalignment between different visual cues.
  \item Better modeling of fingerspelled words and named entities.
  \item Better use of facial and lip information that standard pose-only setups may underuse.
  \item Stronger grounding before the LLM stage, which may reduce generic hallucinated outputs.
\end{itemize}
\end{frame}

\begin{frame}{SignBind-LLM as Our Next Research Step}
\begin{itemize}
  \item Our current work gives us a strong baseline and shows the limitations of simple multimodal alignment.
  \item The next step is not only to scale models, but to improve the architecture itself.
  \item Implementing SignBind-LLM can help us test whether modality-specific experts improve translation quality.
  \item So our future work is moving toward a more structured and high-fidelity multimodal translation system.
\end{itemize}
\end{frame}

\begin{frame}{Conclusion}
\begin{itemize}
  \item Our project evolved from a failed RTMPoses-based pipeline to a better How2Sign-based LLM translation system.
  \item The first stage taught us what was not working.
  \item The second stage gave better metrics, cleaner outputs, and a stronger pipeline.
  \item Our next step is to extend this work with SignBind-LLM style multi-stage modality fusion.
  \item Overall, the project shows both the difficulty and the promise of LLM-based sign language translation.
\end{itemize}
\end{frame}

\begin{frame}{References}
\small
\begin{itemize}
  \item PoseStich-SLT repository: \url{https://github.com/Exploration-Lab/PoseStich-SLT}
  \item How2Sign dataset: \url{https://how2sign.github.io/}
  \item SignBind-LLM paper: \url{https://arxiv.org/abs/2509.00030}
\end{itemize}
\end{frame}

\begin{frame}{Thank You}
\centering
\Large Thank You\\[1em]
\normalsize Questions and Discussion
\end{frame}

\end{document}
